%==============================================================================
%  Beyond the Schoenberg Wall:
%  Signed Bakry--Emery Geometry of Higher-Order Heat Pseudo-Processes
%
%  Paper 3 of the Hurst-line program.
%   Paper 1 (arXiv:2602.19192): fractional-Laplacian curvature = fBM covariance.
%   Paper 2 (tower):            exact graded curvature tower, operator/chaos.
%   Paper 3 (this):             H > 1, where positivity fails -> signed geometry.
%
%  Headline: pseudo-processes. IBM appears only as the first post-wall station.
%
%  Build:  pdflatex hurst_line_signed_curvature.tex   (twice for refs)
%==============================================================================
\documentclass[11pt,reqno]{amsart}

\usepackage[utf8]{inputenc}
\usepackage[T1]{fontenc}
\usepackage{amsmath,amssymb,amsthm,mathtools}
\usepackage{enumitem}
\usepackage[margin=1.15in]{geometry}
\usepackage{xcolor}
\usepackage[hidelinks]{hyperref}

\theoremstyle{plain}
\newtheorem{theorem}{Theorem}[section]
\newtheorem{proposition}[theorem]{Proposition}
\newtheorem{lemma}[theorem]{Lemma}
\newtheorem{corollary}[theorem]{Corollary}
\newtheorem{conjecture}[theorem]{Conjecture}
\theoremstyle{definition}
\newtheorem{definition}[theorem]{Definition}
\newtheorem{remark}[theorem]{Remark}

% honest status box (package-free)
\newenvironment{ledger}
  {\par\medskip\noindent\rule{\linewidth}{0.4pt}\par\nopagebreak\smallskip\small}
  {\par\nopagebreak\smallskip\noindent\rule{\linewidth}{0.4pt}\par\medskip}

\newcommand{\T}{\mathbb{T}}
\newcommand{\R}{\mathbb{R}}
\newcommand{\N}{\mathbb{N}}
\newcommand{\Z}{\mathbb{Z}}
\newcommand{\C}{\mathbb{C}}
\newcommand{\E}{\mathbb{E}}
\newcommand{\HH}{\mathcal{H}}
\newcommand{\KK}{\mathcal{K}}
\newcommand{\Had}{\mathbin{\circ}}
\newcommand{\LP}{\mathcal{LP}}
\newcommand{\Gam}{\Gamma}
\DeclareMathOperator{\spec}{spec}
\DeclareMathOperator{\tr}{tr}
\DeclareMathOperator{\sig}{sig}
\DeclareMathOperator{\rank}{rank}
\DeclareMathOperator{\diag}{diag}

\title[Beyond the Schoenberg wall]{Beyond the Schoenberg Wall:\\
Signed Bakry--\'Emery Geometry of\\ Higher-Order Heat Pseudo-Processes}
\author{Ramiro Fontes}
\address{Quijotic Research}
\email{ramirofontes@gmail.com}
\urladdr{https://www.quijotic.com}
\date{\today}

\subjclass[2020]{Primary 60J76, 47B50; Secondary 60G22, 30C15, 42A82}
\keywords{Bakry--\'Emery curvature; carr\'e du champ; higher-order heat
pseudo-processes; fractional Brownian motion; Schoenberg positivity; Krein space;
indefinite inner product; Laguerre--P\'olya class; real-rootedness}

\begin{document}

\begin{abstract}
On the Hurst line $H=\gamma/2\in(0,\infty)$ the carr\'e du champ of the stable
generator of order $\gamma$ is, on positive frequencies, the fractional-Brownian
covariance kernel $R_H$. For $H\in(0,1]$ this kernel is positive: the regime of
genuine jump processes and fractional Brownian motion, where the graded
positive-side Bakry--\'Emery curvature theory lives. At the heat point
$H=1$ the kernel degenerates to the rank-one form $R_1(s,t)=st$. Past it, for
$H>1$, the kernel is \emph{indefinite}, and the objects of the positive theory ---
the energy space, the tensor tower, the curvature pencil --- continue formally
into a Krein space. This paper develops that continuation.

The location of the wall itself is classical: by Schoenberg's theorem $R_H$ is
positive semidefinite iff $H\in(0,1]$, and we recall it only to fix the boundary.
Our contributions lie past it. (i) A \emph{signature classification} of the even
stations $H=m$: $\Psi_{2m}$ has rank $2m-1$ and signature
$(m-1+\mathbf 1_{m\text{ odd}},\,m-1+\mathbf 1_{m\text{ even}})$, the biharmonic
order $H=2$ being the unique even station retaining a single positive direction,
signature $(1,2)$. (ii) A \emph{curvature--dimension dichotomy}: integrated (weak)
Bakry--\'Emery survives for all $H>1$ and yields Poincar\'e, while the strong
(local) theory is obstructed --- by the signed, non-Markov semigroup throughout
$H>1$, and at $H=2$ additionally by an indefinite pencil whose discriminant is the
very one making the curvature operator's spectrum complex. (iii) To show the
signed structure is not merely descriptive, outputs the spectral gap cannot give:
a \emph{local lower curvature bound}
$\Gam(f,f)(x)\ge\lambda_{\min}(R_H^{(N)})\,\|f\|_{L^2}^2$ whose constant is negative
\emph{exactly} past the wall, a positive-index constraint, and --- beyond the
inertia --- the \emph{full pencil spectrum} $\spec(A_k)$ as the structure constants
of the inter-level relation $\Gam_k=\Re\spec(A_k)\,\Gam_{k-1}
+\Im\spec(A_k)\,\Gam^{\perp}_{k-1}$, with a matching no-go proving $\spec(A_k)$ can
never be a heat-flow decay rate. The signed apparatus thus earns its keep as a
\emph{static} curvature calculus --- provably not a dynamical one. We frame the
signed Bakry--\'Emery data intrinsically (inertia together with the pencil
spectrum), and quarantine all conjectural and speculative material to a fenced
ledger.
\end{abstract}

\maketitle

%==============================================================================
\section{Introduction}\label{sec:intro}
%==============================================================================

\subsection*{The Hurst line}
Let $L_\gamma=-(-\Delta)^{\gamma/2}$ on $\T=\R/2\pi\Z$, $\gamma\in(0,2)$, with
carr\'e du champ $\Gam$ and iterated carr\'e du champ
$\Gam_2=\tfrac12(L_\gamma\Gam-2\Gam(\cdot,L_\gamma\cdot))$. On positive-frequency
trigonometric polynomials its kernel is
\begin{equation}\label{eq:psi}
  \Psi_\gamma(\xi,\eta)=\tfrac12\bigl(|\xi|^\gamma+|\eta|^\gamma-|\xi-\eta|^\gamma\bigr)
  =R_H(\xi,\eta),\qquad H=\tfrac\gamma2,
\end{equation}
the covariance of fractional Brownian motion with Hurst parameter $H$; this is the
identity underlying the positive-side fractional Bakry--\'Emery
theory~\cite{FontesPartI} on the range $0<H<1$, where $R_H$ is a bona fide
covariance.

The kernel \eqref{eq:psi}, however, is defined for every $\gamma$, and the natural
parameter is the Hurst exponent $H=\gamma/2\in(0,\infty)$. Three regimes appear:
\begin{itemize}[leftmargin=1.5em]
\item $H\in(0,1)$ --- genuine pure-jump processes and fractional Brownian motion;
  $R_H$ positive; the positive Bakry--\'Emery theory of~\cite{FontesPartI}.
\item $H=1$ ($\gamma=2$) --- the heat point; $R_1(s,t)=st$ collapses to rank one.
\item $H>1$ --- the \emph{higher-order heat pseudo-processes}, generated by
  $\partial_x^{2m}$ at the even stations $H=m$; $R_H$ \emph{indefinite}. There is
  no probability measure here, only a signed pseudo-process in the sense of
  Hochberg~\cite{Hochberg}, and the geometry is governed by an indefinite inner
  product.
\end{itemize}

The companion paper~\cite{FontesPartI} developed the positive-side theory on
$0<H<1$. The present paper asks what remains beyond the endpoint $H=1$, where
positivity fails. The answer is a
\emph{signed} Bakry--\'Emery geometry: the same carr\'e-du-champ tower
$\Gam_k\leftrightarrow R_H^{\Had k}$, now read in a Krein space, where curvature
is no longer a single positive constant but an inertia.

\subsection*{Results}
Write $R_H=[R_H(n,m)]$, $A^{\Had k}$ for the $k$-th Hadamard power, and
$\sig(\cdot)=(p,q)$ for the inertia (numbers of positive and negative eigenvalues).

\begin{itemize}[leftmargin=1.5em]
\item \textbf{The wall} (Theorem~\ref{thm:wall}). $R_H$ is positive semidefinite
  iff $H\in(0,1]$; at $H=1$ it is rank one, $R_1(s,t)=st$; for $H>1$ it is
  indefinite. The positive theory ends exactly at the heat point.
\item \textbf{Signature of the even stations} (Theorem~\ref{thm:sig}). For
  $H=m\in\N$, $\Psi_{2m}$ has rank $2m-1$ and signature
  $\bigl(m-1+\mathbf 1_{m\text{ odd}},\ m-1+\mathbf 1_{m\text{ even}}\bigr)$. The
  biharmonic order $m=2$ is the unique even order past heat with a single positive
  direction, $\sig=(1,2)$.
\item \textbf{The dictionary} (Theorem~\ref{thm:dict}). For even order $2m$ the
  following are equivalent: (a) $\Psi_{2m}\succeq0$; (b) $e^{t\partial_x^{2m}}$
  preserves the Laguerre--P\'olya class $\LP$; (c) $2m=2$.
\item \textbf{Signed curvature} (Section~\ref{sec:signed}). In the Krein energy
  space the realization $\Gam_k(f,f)=[\,\iota_k\mathbf b,\iota_k\mathbf b\,]$
  survives, but the curvature invariant is the inertia of $R_H^{\Had k}$ together
  with the spectrum of the indefinite pencil $(R_H^{\Had k},R_H^{\Had(k-1)})$,
  which may leave $[0,\infty)$.
\end{itemize}

\begin{remark}[Independence and scope]\label{rem:indep}
This paper may be read independently of the positive-side theory~\cite{FontesPartI};
the needed positive-side facts are recalled as used. The iterated-Brownian / biharmonic case
$H=2$ enters only as the first concrete station past the wall; the organizing
object is the family of higher-order heat pseudo-processes, not any single one.
\end{remark}

%==============================================================================
\section{The Schoenberg wall}\label{sec:wall}
%==============================================================================

The location of the wall is classical --- it is Schoenberg's 1938
theorem~\cite{Schoenberg} in the guise of fBM, recorded here only to fix the
boundary of the signed region. The new content of this paper begins in
Section~\ref{sec:sig}.

\begin{theorem}[The positive theory ends at the heat point;
{\rm \cite{Schoenberg}}]\label{thm:wall}
The fractional-Brownian kernel
$R_H(s,t)=\tfrac12\bigl(|s|^{2H}+|t|^{2H}-|s-t|^{2H}\bigr)$ is positive
semidefinite on $\R$ if and only if $H\in(0,1]$. At $H=1$ it degenerates to the
rank-one form $R_1(s,t)=st$. For $H>1$ it is indefinite.
\end{theorem}

\begin{proof}
By Schoenberg's theorem~\cite{Schoenberg}, $-|s-t|^{\alpha}$ is conditionally
negative definite iff $\alpha\in(0,2]$, and $R_H$ is the covariance obtained from
$|s-t|^{2H}$ by the
standard centering; hence $R_H\succeq0$ iff $2H\le2$, i.e.\ $H\le1$. At $H=1$,
$|s-t|^{2}=s^2-2st+t^2$ gives $R_1(s,t)=\tfrac12(s^2+t^2-(s-t)^2)=st$, rank one and
positive semidefinite. For $H>1$ indefiniteness is exhibited already on two points
by the even-station computation of Theorem~\ref{thm:sig} (and, for non-integer
$H>1$, by a direct $2\times2$ minor); the kernel is no longer a covariance.
\end{proof}

The heat point is thus a genuine \emph{wall}: on $0<H<1$ the admissible energies
form a full-dimensional positive cone (on which the whole Hadamard tower
$R_H^{\Had k}$ is positive), which at $H=1$ collapses to the single line spanned by
$s\mapsto s$. Beyond it, for $H>1$, the cone reopens with the wrong signature.

%==============================================================================
\section{Signature of the even stations}\label{sec:sig}
%==============================================================================

The even stations $H=m$ ($\gamma=2m$) are the generators $\partial_x^{2m}$ of the
higher-order heat pseudo-processes, whose multiplier $e^{tu^{2m}}$ drives the
corresponding flow. Their carr\'e-du-champ kernels are polynomial, and their
inertia is computable in closed form.

\begin{remark}[Which form the signature is of]\label{rem:class}
To keep the paper self-contained (cf.\ Remark~\ref{rem:indep}), we state precisely
the object whose inertia is computed. For $f=\sum_{n\ge1}a_n e^{inx}$ a
positive-frequency trigonometric polynomial, the carr\'e du champ is
$\Gam(f,f)(x)=\sum_{n,m\ge1}a_n\bar a_m e^{i(n-m)x}\,\Psi_\gamma(n,m)$, and we
study the Hermitian \emph{positive-frequency form}
$Q_\gamma(\mathbf a)=\mathbf a^{*}[\Psi_\gamma(n,m)]\mathbf a=\Gam(f,f)(0)$ on
$\mathbf a\in\C^{N}$. Theorem~\ref{thm:sig} computes the signature of
$[\Psi_\gamma(n,m)]$, i.e.\ of $Q_\gamma$. A genuine real function on $\T$ mixes
$\pm n$; its carr\'e du champ is governed by the larger Hermitian kernel
$[\Psi_\gamma(n,m)]_{n,m\in\Z\setminus\{0\}}$, which decomposes under
$n\mapsto-n$ into the positive-frequency block and its conjugate, so the
signed geometry of real functions is two commuting copies of $Q_\gamma$ (the
cross-sign couplings $\Psi_\gamma(n,-m)$ vanish only at the Cauchy point $\gamma=1$
and are part of the positive-side analysis). All statements below concern
$Q_\gamma$ on positive frequencies unless noted.
\end{remark}

\begin{theorem}[Even-order signature]\label{thm:sig}
For $\xi,\eta>0$ and $m\ge1$,
\[
  \Psi_{2m}(\xi,\eta)=\tfrac12\bigl(\xi^{2m}+\eta^{2m}-(\xi-\eta)^{2m}\bigr)
  =\sum_{a=1}^{2m-1}\tfrac12\binom{2m}{a}(-1)^{a+1}\,\xi^{a}\eta^{2m-a},
\]
a kernel of rank $2m-1$ supported on the monomials $\xi,\dots,\xi^{2m-1}$. Its
Gram matrix in the coordinates $(\xi,\dots,\xi^{2m-1})$ is the persymmetric
anti-diagonal matrix $M$ with $M_{a,2m-a}=\tfrac12\binom{2m}{a}(-1)^{a+1}$, and
its signature is
\begin{equation}\label{eq:sig}
  \sig\Psi_{2m}=\bigl(\,m-1+\mathbf 1_{m\textup{ odd}}\,,\ \
  m-1+\mathbf 1_{m\textup{ even}}\,\bigr).
\end{equation}
In particular $m=1$ (heat) gives $(1,0)$ (rank one, positive); $m=2$ gives
$(1,2)$; $m=3$ gives $(3,2)$; $m=4$ gives $(3,4)$. The biharmonic order $m=2$ is
the \emph{unique} even order with exactly one positive direction past the wall.
\end{theorem}

\begin{proof}
The expansion of $(\xi-\eta)^{2m}$ and the cancellation of the $a=0,2m$ terms give
the displayed form; since $(-1)^{2m-a}=(-1)^{a}$, the surviving coefficients are
$\tfrac12\binom{2m}{a}(-1)^{a+1}$, all nonzero, so $\rank M=2m-1$ and $M$ is
supported on the anti-diagonal. $M$ is symmetric: $M_{2m-a,a}
=\tfrac12\binom{2m}{2m-a}(-1)^{(2m-a)+1}=\tfrac12\binom{2m}{a}(-1)^{a+1}=M_{a,2m-a}$.
For each pair $a\in\{1,\dots,m-1\}$ the indices $a$ and $2m-a$ span a
$2\times2$ block $\left(\begin{smallmatrix}0&d_a\\ d_a&0\end{smallmatrix}\right)$
with $d_a=M_{a,2m-a}\neq0$, contributing eigenvalues $\pm|d_a|$ --- one positive,
one negative. The central index $a=m$ contributes the diagonal entry
$M_{m,m}=\tfrac12\binom{2m}{m}(-1)^{m+1}$, a single eigenvalue of sign
$(-1)^{m+1}$: positive if $m$ is odd, negative if $m$ is even. Summing the $m-1$
pairs and the center gives \eqref{eq:sig}.
\end{proof}

For $m=2$ the matrix is
$M=\left(\begin{smallmatrix}0&0&2\\0&-3&0\\2&0&0\end{smallmatrix}\right)$ with
eigenvalues $\{2,-2,-3\}$; already on $\{1,2\}$ one has
$\Psi_4=\left(\begin{smallmatrix}1&8\\8&16\end{smallmatrix}\right)$, $\det=-48<0$.

\begin{remark}[The cone collapses to a line]
For $0<H<1$ the carr\'e du champ is an fBM covariance with a full-dimensional
positive cone carrying the whole positive tower; at $H=1$ it collapses to rank
one. At $H=2$ the cone has collapsed further to
the single energy direction $\Psi_4(\xi,\xi)=\xi^4>0$ against a two-dimensional
negative part. The semigroup $e^{t\partial_x^4}$ no longer contracts an energy: it
has one stabilizing and two destabilizing modes. This is the spectral form of the
fact that $e^{t\partial_x^4}$ does not preserve real-rootedness
(Theorem~\ref{thm:dict}).
\end{remark}

%==============================================================================
\section{The total-positivity dictionary}\label{sec:dict}
%==============================================================================

The Laguerre--P\'olya class $\LP$ is the closure, under locally uniform
convergence, of real polynomials with only real zeros; equivalently, the entire
functions arising as limits of such. A linear flow $e^{tT}$ \emph{preserves}
$\LP$ if it maps $\LP$ into itself for all $t\ge0$. The heat semigroup
$e^{t\partial_x^2}$ preserves $\LP$ (P\'olya); the question past the wall is which
higher-order heat flows do.

\begin{theorem}[Classification of the $\LP$-preserving even stations]\label{thm:dict}
For each integer $m\ge1$ the following are equivalent:
\begin{enumerate}[label=\textup{(\alph*)},leftmargin=2.2em]
\item the carr\'e-du-champ kernel $\Psi_{2m}$ is positive semidefinite;
\item the semigroup $e^{t\partial_x^{2m}}$ preserves the Laguerre--P\'olya class;
\item $2m=2$ (the heat semigroup).
\end{enumerate}
\end{theorem}

\begin{proof}
$(a)\Rightarrow(c)$: by Theorem~\ref{thm:sig}, $\sig\Psi_{2m}=(1,0)$ only at $m=1$;
for $m\ge2$ the signature has a negative part, so $\Psi_{2m}\not\succeq0$.
$(c)\Rightarrow(b)$ is P\'olya's theorem. For $(b)\Rightarrow(c)$ we show $(b)$
fails at every $m\ge2$. The operator $e^{t\partial_x^{2m}}=\varphi(\partial_x)$
has symbol $\varphi(u)=e^{tu^{2m}}$; by the P\'olya--Schur master
theorem~\cite{BorceaBranden}, a differential operator $\varphi(\partial_x)$
preserves $\LP$ iff $\varphi\in\LP$, and $e^{tu^{2m}}$ has order $2m>2$ for
$m\ge2$, hence is not in $\LP$. We make the case $m=2$ fully explicit, in the
paper's spirit: take
\[
  p(x)=(x^2-1)(x^2-4)=x^4-5x^2+4\in\LP\quad(\text{roots }\pm1,\pm2).
\]
Since $p''''=24$ and $p^{(8)}=0$, $e^{t\partial_x^4}p=p+24t=x^4-5x^2+(4+24t)$.
Substituting $u=x^2$, the roots are real iff $25-4(4+24t)\ge0$, i.e.\ iff
$t\le\tfrac{3}{32}$; for $t>\tfrac3{32}$ the roots in $u$ are complex, so
$e^{t\partial_x^4}p\notin\LP$. Thus heat ($m=1$) is the last even station to
preserve $\LP$, and the three statements are equivalent to $m=1$.
\end{proof}

\begin{remark}
Theorem~\ref{thm:dict} is a \emph{classification}, not an identification: both
kernel positivity and $\LP$-preservation single out $m=1$, so their failure loci
coincide, but the two conditions are genuinely different objects. Indeed
\S\ref{sec:dichotomy} exhibits energy-positive functions that are not real-rooted,
so the energy cone strictly contains the real-rooted one. What the dictionary says
is the modest and correct thing: among even stations, the heat flow is the last to
keep \emph{either} property.
\end{remark}

%==============================================================================
\section{Signed Bakry--\'Emery geometry}\label{sec:signed}
%==============================================================================

We now state what the curvature tower becomes past the wall. The positive-side
construction is recalled as follows: an energy space $\HH=\R^N$ with
$\langle e_n,e_m\rangle=R_H(n,m)$, a stochastic representative
$\Phi_f(x)=\sum_n b_n(x)e_n$, and the diagonal tensor lift
$\iota_k(\mathbf v)=\sum_n v_n e_n^{\otimes k}$, for which
$\Gam_k(f,f)=\langle\iota_k\mathbf b,\iota_k\mathbf b\rangle_{\HH^{\otimes k}}
=\mathbf b^{*}R_H^{\Had k}\mathbf b$.

For $H>1$ the form $\langle\cdot,\cdot\rangle$ is no longer positive: by
Theorem~\ref{thm:wall}, $R_H$ is indefinite, so $\HH$ is a finite-dimensional
\emph{Krein space} --- a nondegenerate inner-product space with indefinite form
$[\cdot,\cdot]$ and fundamental symmetry $J$ ($J=J^{*}$, $J^{2}=I$,
$[x,y]=\langle Jx,y\rangle$) \cite{Bognar}. The tensor powers $\HH^{\otimes k}$
inherit indefinite forms with Gram matrices $R_H^{\Had k}$, and the realization
\begin{equation}\label{eq:krein-real}
  \Gam_k(f,f)(x)=[\,\iota_k\mathbf b(x),\,\iota_k\mathbf b(x)\,]_{\HH^{\otimes k}}
  =\mathbf b(x)^{*}R_H^{\Had k}\,\mathbf b(x)
\end{equation}
survives verbatim --- it never used positivity, only the bilinear pairing. What
does not survive is the \emph{interpretation} of curvature, and the purpose of
this section is to say precisely what replaces it.

\subsection{The curvature pencil is $J$-self-adjoint}
The level-$k$ curvature on the positive range was the optimal constant in
$\Gam_k\ge\kappa\,\Gam_{k-1}$, i.e.\ the smallest generalized eigenvalue of the
pencil $(R_H^{\Had k},R_H^{\Had(k-1)})$. Past the wall the same pencil defines an
operator, but its spectrum is no longer real.

\begin{proposition}[Spectral symmetry of the curvature operator]\label{prop:Jself}
Assume $R_H^{\Had(k-1)}$ is nonsingular and let
$A_k:=\bigl(R_H^{\Had(k-1)}\bigr)^{-1}R_H^{\Had k}$. Then $A_k$ is self-adjoint
with respect to the indefinite form $G:=R_H^{\Had(k-1)}$, i.e.\
$G A_k=A_k^{\top}G$ (so $A_k$ is $J$-self-adjoint in the Krein space with Gram
$G$). Consequently $\spec(A_k)$ is symmetric under complex conjugation: nonreal
eigenvalues occur in conjugate pairs $\lambda,\bar\lambda$, and real eigenvalues
may have either sign.
\end{proposition}

\begin{proof}
$GA_k=R_H^{\Had k}$ is symmetric, and $A_k^{\top}G=R_H^{\Had k}\,
\bigl(R_H^{\Had(k-1)}\bigr)^{-1}R_H^{\Had(k-1)}=R_H^{\Had k}$; hence
$GA_k=A_k^{\top}G$. A $G$-self-adjoint operator with real $G$ has a real
characteristic polynomial, so its spectrum is conjugation-symmetric; for definite
$G$ this forces real spectrum, but for indefinite $G$ conjugate complex pairs are
admissible.
\end{proof}

\begin{definition}[Signed Bakry--\'Emery data]\label{def:signed}
For $H>1$ the signed curvature data of the tower at level $k$ is the triple
\[
  \Bigl(\ \sig R_H^{\Had k},\quad \spec(A_k),\quad \kappa_+^{(k)}\ \Bigr),
\]
consisting of the inertia of the level-$k$ form, the spectrum of the
$J$-self-adjoint curvature operator $A_k$ of Proposition~\ref{prop:Jself}, and the
\emph{positive-cone curvature}
\[
  \kappa_+^{(k)}:=\inf\Bigl\{\ \tfrac{\mathbf v^{*}R_H^{\Had k}\mathbf v}
  {\mathbf v^{*}R_H^{\Had(k-1)}\mathbf v}\ :\
  \mathbf v^{*}R_H^{\Had(k-1)}\mathbf v>0\ \Bigr\},
\]
the best curvature constant on the maximal positive subspace of $\Gam_{k-1}$.
\end{definition}

We stress that Definition~\ref{def:signed} is \emph{proposed data}, not yet a
curvature--dimension theory. In particular $\kappa_+^{(k)}$ need not be finite:
as $\mathbf v$ approaches the boundary of the positive cone, the denominator
$\mathbf v^{*}R_H^{\Had(k-1)}\mathbf v\downarrow0$ while the numerator stays
bounded away from zero, so the infimum is generically $-\infty$. Thus
$\kappa_+^{(k)}$ is canonical as a definition but is not, on its own, a usable
curvature lower bound; whether some renormalization of it controls a functional
inequality is exactly the open problem (\S\ref{sec:dichotomy}).

On the positive range all three pieces collapse to the single number
$\kappa^{(k)}=\lambda_{\min}>0$; past the wall they genuinely separate, and the
following computation shows that the spectrum of $A_k$ is the part that fails to
be real.

\begin{proposition}[The biharmonic curvature is complex]\label{prop:complex}
At the biharmonic station $H=2$, on the two frequencies $\{1,2\}$, the level-two
carr\'e-du-champ data are
\[
  R_2=\Psi_4=\begin{pmatrix}1&8\\8&16\end{pmatrix},\qquad
  \Psi_4^{\Had2}=\begin{pmatrix}1&64\\64&256\end{pmatrix},
\]
and the curvature operator $A_2=\Psi_4^{-1}\Psi_4^{\Had2}$ has
$\tr A_2=\tfrac{47}{3}$, $\det A_2=80$, hence eigenvalues
\[
  \spec(A_2)=\Bigl\{\ \tfrac{47\pm i\sqrt{671}}{6}\ \Bigr\}
  \;\approx\;7.83\pm 4.32\,i\ \notin\ \R .
\]
In particular there is no real curvature constant at $H=2$: the curvature operator
of the first crossing of the wall is genuinely complex.
\end{proposition}

\begin{proof}
$\det\Psi_4=-48$ and $\det\Psi_4^{\Had2}=-3840$, so $\det A_2=\det\Psi_4^{\Had2}/
\det\Psi_4=80$; a direct multiplication gives
$A_2=\tfrac1{-48}\left(\begin{smallmatrix}-496&-1024\\56&-256\end{smallmatrix}\right)$
with $\tr A_2=752/48=47/3$. The characteristic polynomial
$\lambda^{2}-\tfrac{47}{3}\lambda+80$ has discriminant $-\tfrac{671}{9}<0$, giving
the stated complex-conjugate pair.
\end{proof}

\begin{ledger}
\noindent\textbf{The central open problem.}
On the positive range, $\mathrm{CD}(\kappa,\infty)$ drives the Bakry--\'Emery
machine: a curvature lower bound yields gradient contraction, Poincar\'e, and
log-Sobolev inequalities for the semigroup \cite{BGL}. Past the wall that machine
has no input: by Proposition~\ref{prop:complex} the curvature operator $A_k$ is
$J$-self-adjoint with genuinely complex spectrum, so there is no $\kappa$ to feed
it. Three candidate replacements are on the table --- the inertia
\eqref{eq:sig}, the spectral radius / real-part envelope of $\spec(A_k)$, and the
positive-cone curvature $\kappa_+^{(k)}$ of Definition~\ref{def:signed} --- and the
problem this paper leaves open is to determine which (if any) controls a
semigroup functional inequality for the signed pseudo-process. Turning
Definition~\ref{def:signed} into a curvature--dimension calculus for Krein-space
Dirichlet forms is the natural next theorem, not a corollary of the present one.
\end{ledger}

\begin{remark}[The mildness of the biharmonic crossing]
Among the even stations, $m=2$ is distinguished twice: it is the unique order with
a single positive direction (Theorem~\ref{thm:sig}), and it is the first station
past heat, so its negative part is only two-dimensional. In this precise sense the
iterated-Brownian / biharmonic pseudo-process is the \emph{mildest} crossing of
the wall --- one positive direction still standing --- but, as
\eqref{eq:sig} shows, this mildness is special to it and not the beginning of a
pattern: at $m\ge3$ the positive and negative parts are nearly balanced.
\end{remark}

%==============================================================================
\section{The Krein--Bakry--\'Emery dichotomy}\label{sec:dichotomy}
%==============================================================================

Section~\ref{sec:signed} left curvature past the wall as \emph{proposed data}.
We now determine what part of the Bakry--\'Emery machine actually survives. The
outcome is a sharp split: the \emph{integrated} (weak) theory survives the wall
unconditionally for all $H>1$, while the \emph{local} (strong) theory --- the half
that yields log-Sobolev and hypercontractivity --- is obstructed for all $H>1$ by
the signed semigroup, and at the biharmonic station $H=2$ by the curvature
signature as well.

\subsection{Two conditions}
Integrating the carr\'e-du-champ identities over $\T$ collapses the off-diagonal
of $R_H^{\Had k}$, where the indefiniteness lives, and leaves the manifestly
positive quantities
\[
  \mathcal E_1(f)=\int_\T\Gam(f,f)\,d\mu=\langle f,-Lf\rangle=\sum_n\lambda_n|\hat f_n|^2,
  \qquad
  \mathcal E_2(f)=\int_\T\Gam_2(f,f)\,d\mu=\|Lf\|^2=\sum_n\lambda_n^2|\hat f_n|^2,
\]
where $-Le_n=\lambda_n e_n$, $\lambda_n=|n|^{\gamma}\ge1$. Define
\[
  \mathrm{CD}_{\mathrm{int}}(\kappa):\ \mathcal E_2\ge\kappa\,\mathcal E_1\ \forall f,
  \qquad
  \mathrm{CD}_{\mathrm{loc}}(\kappa):\ \Gam_2(f,f)(x)\ge\kappa\,\Gam(f,f)(x)\ \forall f,x
  \ \ \Longleftrightarrow\ \ R_H^{\Had2}\succeq\kappa R_H .
\]

\begin{theorem}[Dichotomy at the wall]\label{thm:dichotomy}
\textup{(i)} $\mathrm{CD}_{\mathrm{int}}(1)$ holds for \emph{all} $H>0$ --- it is
the spectral gap $\lambda_n^2\ge\lambda_1\lambda_n$ with $\lambda_1=1$ --- and
yields the Poincar\'e inequality $\mathrm{Var}(f)\le\mathcal E_1(f)$ for the
pseudo-process, unconditionally past the wall.
\textup{(ii)} At the biharmonic station $H=2$, $\mathrm{CD}_{\mathrm{loc}}(\kappa)$
fails for every real $\kappa$.
\end{theorem}

\begin{proof}
(i) is immediate from $\lambda_n\ge\lambda_1=1$. For (ii), the principal
$\{1,2\}$ minor requires $R^{\Had2}-\kappa R\succeq0$, hence
$\det\ge0$: $(1-\kappa)(256-16\kappa)-(64-8\kappa)^2\ge0$, i.e.\
$\kappa^2-\tfrac{47}{3}\kappa+80\le0$, whose discriminant $-\tfrac{671}{9}<0$
makes it unsatisfiable; this minor obstructs $R^{\Had2}-\kappa R\succeq0$ for all
$N\ge2$.
\end{proof}

\begin{remark}[The uniform statement over $H>1$ is open --- and subtle near the wall]
\label{rem:uniform}
We do \emph{not} claim $\mathrm{CD}_{\mathrm{loc}}$ fails throughout the signed
region: it does not follow from the $\{1,2\}$ minor. With $a:=2^{2H}$ that minor
is $R_H=\bigl(\begin{smallmatrix}1&a/2\\a/2&a\end{smallmatrix}\bigr)$,
$R_H^{\Had2}=\bigl(\begin{smallmatrix}1&a^2/4\\a^2/4&a^2\end{smallmatrix}\bigr)$,
and $\det(R_H^{\Had2}-\kappa R_H)$ is a downward parabola in $\kappa$ (for $a>4$)
with discriminant proportional to $-(a^3-6a^2+8a-4)$. That cubic is negative for
$a\in(4,a_*)$, $a_*\approx4.38$, i.e.\ for $H\in(1,\ {\approx}\,1.065)$, where the
parabola has real roots and the minor \emph{admits} a real $\kappa$ (e.g.\ at
$H=1.03$ every $\kappa\in[-19.1,-0.45]$ makes the $\{1,2\}$ block PSD). It first
obstructs at $a=a_*$ and is obstructed at $H=2$ ($a=16$). So whether the
\emph{form} obstruction holds throughout $H>1$ --- or the strong theory survives,
with negative curvature, in a layer just above the wall --- requires a separate
finite-minor argument and is genuinely open near $H=1^+$. The dichotomy below does
not depend on it: the strong theory is killed for \emph{all} $H>1$ by the second
obstruction (O2), independently of any curvature bound.
\end{remark}

\begin{proposition}[The strong-side obstruction is the complex spectrum]\label{prop:disc}
The discriminant $-\tfrac{671}{9}$ of Theorem~\ref{thm:dichotomy}(ii) is exactly
the one making $\spec(A_2)=\tfrac{47\pm i\sqrt{671}}{6}$ non-real
(Proposition~\ref{prop:complex}): on a $2\times2$ block, a real local curvature
constant exists iff the curvature pencil there has real spectrum. The biharmonic
block has complex spectrum, so no real local constant exists --- the obstruction
of Theorem~\ref{thm:dichotomy}(ii) and the complex eigenvalues of
Proposition~\ref{prop:complex} are one computation.
\end{proposition}

\subsection{Why the strong theory fails}
The strong Bakry--\'Emery conclusions (gradient contraction, log-Sobolev) come
from the interpolation
$\partial_s\bigl[P_s\,\Gam(P_{t-s}f)\bigr]=2\,P_s\,\Gam_2(P_{t-s}f)$, which needs
\emph{two} things: a pointwise bound $\Gam_2\ge\kappa\Gam$, and a
positivity-preserving $P_s$ to push it through. Past the wall the second fails
universally.

\smallskip\noindent\textbf{(O2) Non-Markov semigroup --- for all $H>1$.} For
$\gamma>2$ the kernel of $e^{tL}=e^{-t(-\Delta)^{\gamma/2}}$ is \emph{signed}
(this is precisely the pseudo-process), so $e^{tL}$ is not positivity-preserving.
A nonnegative integrand need not stay nonnegative under the signed $P_s$, so the
interpolation breaks \emph{even if} a local curvature bound held. This obstruction
holds throughout the signed region, independently of any form inequality, and is
what makes the strong theory fail for every $H>1$.

\smallskip\noindent\textbf{(O1) Indefinite form --- at the biharmonic station.}
At $H=2$ there is in addition no curvature bound to feed the interpolation:
by Proposition~\ref{prop:disc}, $\mathrm{CD}_{\mathrm{loc}}(\kappa)$ holds for no
real $\kappa$. Whether (O1) extends to all $H>1$ is open
(Remark~\ref{rem:uniform}); (O2) does not need it.

\smallskip
The integrated identity $\tfrac{d}{ds}\mathcal E_1(P_sf)=-2\,\mathcal E_2(P_sf)$
needs neither obstruction and powers the surviving Poincar\'e decay. So past the
wall the weak Bakry--\'Emery theory is intact for all $H>1$, while the strong
theory (log-Sobolev, hypercontractivity) is obstructed for all $H>1$ by (O2) ---
and at $H=2$ doubly so, by (O1) as well.

\subsection{What the signed structure produces}\label{sub:produces}
The surviving Poincar\'e inequality of Theorem~\ref{thm:dichotomy}(i) uses none of
the signed apparatus: it is the spectral gap $\lambda_n\ge\lambda_1=1$, true for
every $H>0$ and provable without ever mentioning inertia or the indefinite pencil.
It is fair to ask what, if anything, the signed geometry adds. The spectral gap
controls only the \emph{mean} $\int_\T\Gam(f,f)\,d\mu=\mathcal E_1\ge0$; it is
blind to the carr\'e du champ \emph{as a function of $x$}, which dips negative past
the wall. That pointwise behaviour is exactly what the inertia governs.

\begin{proposition}[Two outputs of the signed structure]\label{prop:produces}
Let $f=\sum_{n=1}^{N}a_ne^{inx}$, so $\Gam(f,f)(x)=\mathbf b(x)^{*}R_H\mathbf b(x)$
with $\mathbf b(x)_n=a_ne^{inx}$ and $\|\mathbf b(x)\|^2=\sum_n|a_n|^2=\|f\|^2_{L^2}$.
\begin{enumerate}[label=\textup{(\roman*)},leftmargin=2.2em]
\item \textbf{Local lower curvature bound.} For every $x$,
\begin{equation}\label{eq:floor}
  \Gam(f,f)(x)\ \ge\ \lambda_{\min}\!\bigl(R_H^{(N)}\bigr)\,\|f\|^2_{L^2},
\end{equation}
and the constant is optimal. Its sign \emph{detects the wall}:
$\lambda_{\min}(R_H^{(N)})\ge0$ for $H\le1$ \textup(so \eqref{eq:floor} is the
vacuous $\Gam\ge0$ there\textup), while $\lambda_{\min}(R_H^{(N)})<0$ for $H>1$,
quantifying the negative excursion of the carr\'e du champ. The spectral gap is
silent about this floor, as it controls only $\int\Gam\,d\mu$.
\item \textbf{Positive-index constraint.} At $H=2$, the form $Q_4$ has positive
index $1$ (Theorem~\ref{thm:sig}). Hence for any two linearly independent
$f,g\in\mathcal T_N^{+}$ there is a nontrivial combination $h=\alpha f+\beta g$
with $\Gam(h,h)(0)\le0$. In the positive theory $H<1$ no such $h$ exists
($\Gam\ge0$ always); the constraint is a genuine feature of the signed regime.
\end{enumerate}
\end{proposition}

\begin{proof}
(i) is the Rayleigh bound $\mathbf b^{*}R_H\mathbf b\ge\lambda_{\min}(R_H)
\|\mathbf b\|^2$, with equality when $\mathbf a$ is a bottom eigenvector and $x=0$;
the sign statement is Theorem~\ref{thm:wall} together with
$\lambda_{\min}(R_H)<0\Leftrightarrow R_H\not\succeq0$. (ii) is the Courant--Fischer
index theorem: if the positive eigenspace of $R_H$ has dimension $1$, every
$2$-dimensional subspace meets the non-positive part nontrivially; take
$\mathbf a=\alpha\hat f+\beta\hat g$ there and read $\Gam(h,h)(0)=\mathbf a^{*}
R_H\mathbf a\le0$.
\end{proof}

These are modest --- a lower bound and an index count --- but they are the first
statements in the paper that the spectral gap cannot reach, and their constants
are inertia data of the signed form. They use only the bottom eigenvalue of
$R_H$; the \emph{full} pencil spectrum does more. Section~\ref{sub:pencil} shows
that $\spec(A_k)$ --- real and imaginary parts both --- is the set of structure
constants relating consecutive carr\'e-du-champ levels, while a matching no-go
(Proposition~\ref{prop:nogo}) shows it can nonetheless never be a decay rate of the
heat flow: the vindication is static, and provably not dynamical.

\subsection{The pencil spectrum as static curvature data}\label{sub:pencil}
The inertia outputs of \S\ref{sub:produces} use only the extreme eigenvalues of
$R_H$. The full pencil spectrum --- the object that distinguishes the signed
theory --- turns out to be the set of structure constants tying $\Gam_k$ to
$\Gam_{k-1}$, with its real and imaginary parts both carrying meaning. Write
$[\mathbf x,\mathbf y]_{k-1}=\mathbf x^{*}R_H^{\Had(k-1)}\mathbf y$ and
$\Gam_j(f,f)(0)=\mathbf b^{*}R_H^{\Had j}\mathbf b$, $\mathbf b=(a_n)$.

\begin{proposition}[Spectral decomposition of the curvature tower]\label{prop:decomp}
Let $W$ be a $2$-dimensional $[\cdot,\cdot]_{k-1}$-nondegenerate, $A_k$-invariant
subspace on which $A_k=(R_H^{\Had(k-1)})^{-1}R_H^{\Had k}$ has the complex
eigenpair $\rho_k\pm i r_k$ \textup($r_k>0$\textup). On $W$:
\begin{enumerate}[label=\textup{(\roman*)},leftmargin=2.2em]
\item $A_k=\rho_k I+r_k J_0$ with $J_0$ \emph{$[\cdot,\cdot]_{k-1}$-self-adjoint},
$J_0^2=-I$, and $[J_0\mathbf b,J_0\mathbf b]_{k-1}=-[\mathbf b,\mathbf b]_{k-1}$:
the complex structure $J_0$ \emph{reverses the sign} of the carr\'e du champ.
\item With the conjugate energy
$\Gam^{\perp}_{k-1}(f):=[\mathbf b,J_0\mathbf b]_{k-1}$,
\begin{equation}\label{eq:decomp}
  \Gam_k(f,f)(0)=\rho_k\,\Gam_{k-1}(f,f)(0)+r_k\,\Gam^{\perp}_{k-1}(f),
\end{equation}
so $\Re\spec(A_k)$ and $\Im\spec(A_k)$ are exactly the coefficients of the
carr\'e du champ and of its conjugate.
\item The pair obeys the sharp curvature ellipse
\begin{equation}\label{eq:ellipse}
  \bigl(\Gam_k-\rho_k\Gam_{k-1}\bigr)^2+r_k^2\,\Gam_{k-1}^2
  = r_k^2\,\bigl|\det R_H^{\Had(k-1)}\big|_W\bigr|\,\mathrm{area}_W(\mathbf b)^2,
\end{equation}
centred at $\rho_k\Gam_{k-1}$. At $H=2$ on $\{1,2\}$:
$\rho_2=\tfrac{47}{6}$, $r_2=\tfrac{\sqrt{671}}{6}$, $|\det R|=48$.
\end{enumerate}
\end{proposition}

\begin{proof}
(i) Cayley--Hamilton for the real $2\times2$ block $A_k|_W$ with eigenvalues
$\rho_k\pm ir_k$ gives $(A_k-\rho_k)^2=-r_k^2 I$, so $J_0:=(A_k-\rho_k)/r_k$ has
$J_0^2=-I$; it is $[\cdot,\cdot]_{k-1}$-self-adjoint because $A_k$ is
(Proposition~\ref{prop:Jself}) and $\rho_k,r_k$ are real, whence
$[J_0\mathbf b,J_0\mathbf b]_{k-1}=[\mathbf b,J_0^2\mathbf b]_{k-1}
=-[\mathbf b,\mathbf b]_{k-1}$. (ii) $\Gam_k(f,f)(0)=\mathbf b^{*}R_H^{\Had k}\mathbf b
=[\mathbf b,A_k\mathbf b]_{k-1}=\rho_k[\mathbf b,\mathbf b]_{k-1}
+r_k[\mathbf b,J_0\mathbf b]_{k-1}$. (iii) In the basis $\{\mathbf b,J_0\mathbf b\}$
the form $[\cdot,\cdot]_{k-1}$ has Gram
$\left(\begin{smallmatrix}\Gam_{k-1}&\Gam^{\perp}\\ \Gam^{\perp}&-\Gam_{k-1}\end{smallmatrix}\right)$,
whose determinant equals
$\det(R_H^{\Had(k-1)}|_W)\,|\det[\mathbf b\,|\,J_0\mathbf b]|^2$ by change of basis;
as $W$ carries a complex pair it is hyperbolic, $\det(R_H^{\Had(k-1)}|_W)<0$, and
$-\Gam_{k-1}^2-(\Gam^{\perp})^2=\det(R_H^{\Had(k-1)}|_W)\,\mathrm{area}_W^2$
rearranges to the circle identity; substituting
$\Gam^{\perp}=(\Gam_k-\rho_k\Gam_{k-1})/r_k$ gives \eqref{eq:ellipse}.
\end{proof}

\begin{proposition}[No heat-flow rate equals $\spec(A_k)$]\label{prop:nogo}
The semigroup $P_t=e^{tL}$ is diagonal, $\widehat{P_tf}_n=e^{-\lambda_n t}\hat f_n$.
Hence every integrated energy
$\mathcal E_j(P_tf)=\sum_n\lambda_n^{j}e^{-2\lambda_n t}|\hat f_n|^2$ and every
pointwise level $\Gam_k(P_tf,P_tf)(x)
=\sum_{n,m}e^{-(\lambda_n+\lambda_m)t}\hat f_n\bar{\hat f}_m e^{i(n-m)x}R_H^{\Had k}(n,m)$
decay through the generator rates $\{2\lambda_n\}$ and $\{\lambda_n+\lambda_m\}$
only. Since $\spec(A_k)$ is built from Hadamard ratios of the $\Psi_\gamma(n,m)$ ---
which enter the \emph{coefficients}, never the exponents --- no such functional
decays at a rate equal to $\Re\spec(A_k)$. In particular the Poincar\'e/decay rate
of $\mathcal E_1$ is the spectral gap $2\lambda_1$, and no Krein--Poincar\'e or
gradient inequality for $P_t$ can carry the constant $\Re\spec(A_k)$.
\end{proposition}

\begin{remark}[Vindication, exactly bounded]
Together, Propositions~\ref{prop:decomp} and~\ref{prop:nogo} settle the status of
the indefinite pencil. \emph{Statically} it is the genuine curvature datum past the
wall: its full complex spectrum is the set of structure constants of the
inter-level relation \eqref{eq:decomp}, with $\Re\spec(A_k)$ the centre and
$\Im\spec(A_k)$ the amplitude of the conjugate energy the wall switches on ---
not the inertia of \S\ref{sub:produces}, but the whole eigenvalue. \emph{Dynamically}
it does no work: the only flow on which $\spec(A_k)$ is a rate is the internal
Krein rotation $e^{sJ_0}$, under which $(\Gam_{k-1},\Gam^{\perp}_{k-1})$ revolves on
the circle \eqref{eq:ellipse} and $\Gam_k$ oscillates about $\rho_k\Gam_{k-1}$ with
amplitude $r_k|\Gam^{\perp}_{k-1}|$. The pencil governs the curvature's internal
geometry, not the heat semigroup's decay --- which is the structural reason the
strong theory does not return, and the precise sense in which the signed apparatus
earns its keep: as a static curvature calculus, not a dynamical one.
\end{remark}

\subsection{Curvature versus real-rootedness}
Real-rootedness is governed by a different cone. Write
$\mathcal C=\{f:\Gam(f,f)(x)>0\ \forall x\}$ for the energy-positive cone and
$\LP$ for the real-rooted (Laguerre--P\'olya) cone.

\begin{proposition}[Energy-positivity is strictly weaker than real-rootedness]
\label{prop:LPC}
At $H=2$ on the frequencies $\{1,2\}$, with $f=a_1e^{ix}+a_2e^{2ix}$,
\[
  \Gam(f,f)(x)=|a_1|^2+16|a_2|^2+16\,\Re\!\bigl(a_1\bar a_2e^{-ix}\bigr).
\]
The zero of $f=e^{ix}(a_1+a_2e^{ix})$ sits at $e^{ix}=-a_1/a_2$, which corresponds
to a real $x$ iff it lies on the unit circle, i.e.\ iff $|a_1|=|a_2|$. Every such
real-rooted $f$ has $\min_x\Gam(f,f)=|a_1|^2>0$, so $\LP\subseteq\mathcal C$. But
$f_0=e^{ix}+2e^{2ix}$ has $\Gam(f_0,f_0)(x)=65+32\cos x\ge33>0$, so
$f_0\in\mathcal C$, while $|a_1|=1\ne2=|a_2|$ puts its zero off the unit circle
($x\notin\R$), so $f_0\notin\LP$. Hence $\LP\subsetneq\mathcal C$.
\end{proposition}

So the curvature tower certifies only membership in $\mathcal C$ --- a condition
\emph{necessary but not sufficient} for real-rootedness. It cannot tell $f_0$
apart from a real-rooted competitor. The gap $\mathcal C\setminus\LP$ is exactly
the test-function dependence the inertia does not see. In particular, for any
threshold programme based on a flow $P_t$, the cone-exit time
$\Lambda_\Gamma(f)=\inf\{T:P_tf\in\mathcal C\ \forall t\ge T\}$ is at most the
real-rootedness threshold; so curvature-positivity can furnish a \emph{necessary}
condition for such problems --- the de Bruijn--Newman setting among them
(\S\ref{sec:ledger}) --- but never a sufficient one.

\subsection{The negative-dimension escape, and its limit}
A dimension term $\tfrac1n(Lf)^2$ with $n>0$ \emph{strengthens}
$\mathrm{CD}(\kappa,\infty)$ and cannot help. Negative effective dimension $n<0$
weakens it: with $\lambda=(\lambda_n)_n$ ($\lambda_n=n^4$ at $H=2$), the dimension
term has rank-one Gram $\lambda\lambda^\top$, and $\mathrm{CD}(\kappa,n)$ becomes
\[
  R_H^{\Had2}-\kappa R_H+\tfrac1{|n|}\,\lambda\lambda^\top\ \succeq\ 0 .
\]

\begin{proposition}[A rank-one repair, of bounded reach]\label{prop:dim}
\textup{(i)} At $N=2$, $H=2$, the choice $\kappa=-100$, $|n|^{-1}=9000$
\textup(i.e.\ $n=-\tfrac1{9000}$\textup) gives a matrix with positive diagonal and
$\det=16960>0$, so $\mathrm{CD}(\kappa,n)$ holds: the biharmonic crossing admits a
finite negative-dimension curvature bound.
\textup{(ii)} The correction $\tfrac1{|n|}\lambda\lambda^\top$ is rank one, so it
raises the inertia by at most one negative direction; since the negative inertia
of $R_H^{\Had2}-\kappa R_H$ grows with $N$, for $N$ large \emph{no} finite $n$
repairs $\mathrm{CD}_{\mathrm{loc}}$.
\textup{(iii)} Where it does hold, $\kappa<0$, so it yields at best a
Brascamp--Lieb / weighted-Poincar\'e inequality, never log-Sobolev; and it leaves
\textup{(O2)} untouched.
\end{proposition}

The moral is structural: a \emph{scalar} dimension, even a negative one, is
exactly one rank short of the problem. Repairing the growing negative inertia
would require a matrix-valued curvature--dimension condition
$\mathrm{CD}(\kappa,\mathbf n)$, and even then the non-Markov obstruction (O2)
would keep the strong inequalities out of reach. This is the precise sense in
which the open problem of Section~\ref{sec:signed} is settled: there \emph{is} a
Bakry--\'Emery theory past the wall, it is the weak (integrated) theory; the
strong theory is obstructed for all $H>1$ by the signed semigroup (O2), and at the
biharmonic station additionally by the curvature signature (O1), in a way no
scalar curvature--dimension bound can repair. (Whether the form obstruction (O1)
extends throughout $H>1$, or the strong theory partially survives just above the
wall, is left open by Remark~\ref{rem:uniform}.)

%==============================================================================
\section{Ledger: proved, conjectured, dreamt}\label{sec:ledger}
%==============================================================================

The grand motivation for crossing the wall is the de Bruijn--Newman circle of
questions: the even stations $H=m$ are exactly the higher-order analogues of the
heat flow whose $H=1$ instance carries the de Bruijn--Newman constant $\Lambda$
(known to be nonnegative~\cite{RodgersTao}), and one would like to read the
location of zeros off the signed curvature. We are
deliberately conservative about how far this is established.

\begin{ledger}
\noindent\textbf{Proved here} (and independent of any speculation):
the Schoenberg wall (Theorem~\ref{thm:wall}); the even-order signature
\eqref{eq:sig} (Theorem~\ref{thm:sig}); the total-positivity dictionary
(Theorem~\ref{thm:dict}); the Krein-space realization \eqref{eq:krein-real} of the
carr\'e-du-champ tower past the wall; and the dichotomy of
Section~\ref{sec:dichotomy} --- weak (integrated) Bakry--\'Emery survives the wall
and gives Poincar\'e for all $H>1$ (Theorem~\ref{thm:dichotomy}), the strong
(local) theory is killed for all $H>1$ by the non-Markov semigroup (O2) and, at
the biharmonic station $H=2$, additionally by the signature/complex-spectrum
obstruction (Proposition~\ref{prop:disc}), the energy cone strictly contains the
real-rooted one (Proposition~\ref{prop:LPC}), and the negative-dimension repair is
rank-one, local to small $N$, and cannot reach log-Sobolev
(Proposition~\ref{prop:dim}); and the static/dynamical resolution of the pencil ---
$\spec(A_k)$ is the structure constants of the inter-level curvature relation
(Proposition~\ref{prop:decomp}) but provably never a heat-flow rate
(Proposition~\ref{prop:nogo}).

\smallskip\noindent\textbf{Conjectured} (computationally supported, not proved):
that $\gamma=1$ maximizes the level-$k$ curvature for every $k$ on the positive
range (the level-$2$ case is the phase transition of~\cite{FontesPartI}); and that
the higher-order de Bruijn--Newman constants $\Lambda(2m)$ are finite, with the
quartic value $\Lambda(4)$ governed by the detachment of the saddle of
$g(u)=tu^{4}+\log\Phi(u)$ at $u_*=\tfrac34$, $t_*=\tfrac{64}{27}\pi e^{3}$. We
flag explicitly that the first form of the associated $\Lambda(4)$
``re-coherence'' law was tested and \emph{refuted}; what is supported is finiteness
via the saddle mechanism, with uniformity in $t$ unproved.

\smallskip\noindent\textbf{Speculative} (the frontier, recorded as such):
the reading of the Riemann hypothesis as a critical tangency of the curvature
flow; zero dynamics as a curvature/Ricci flow on the Hurst line; and the existence
of a single analytic operator interpolating the positive and signed sides across
the wall $H=1$. The earlier claim that ``tower positivity \emph{is}
real-rootedness'' is \emph{false} --- tower positivity is test-function
independent --- and has been replaced by the dictionary of
Section~\ref{sec:dict} and the positive-side critical-line reading (the secular
polynomial $z^{k}-(z-1)^{k}$, all of whose zeros lie on $\Re=\tfrac12$).
\end{ledger}

\begin{remark}[Outlook]
The signed structure has earned a definite, if bounded, role. It yields
inequalities beyond the spectral gap (Proposition~\ref{prop:produces}), and its
pencil spectrum is the static curvature datum of the tower
(Proposition~\ref{prop:decomp}) --- while a no-go (Proposition~\ref{prop:nogo})
shows it cannot be a heat-flow rate, so the calculus is static, not dynamical. Two
questions remain genuinely open. First, whether the local obstruction holds
throughout $H>1$ or the strong theory survives in a layer just above the wall
(Remark~\ref{rem:uniform}). Second --- the natural way to seek a dynamical theory ---
whether the internal Krein rotation $e^{sJ_0}$ can be coupled to the heat flow on a
doubled (complexified) state space, so that $\spec(A_k)$ becomes a genuine rate of
some semigroup; the no-go forbids this for $P_t$ alone, not for an enlarged
dynamics.
\end{remark}

%==============================================================================
\begin{thebibliography}{9}

\bibitem{FontesPartI}
R.~Fontes,
\emph{Bakry--\'Emery curvature of the fractional Laplacian via fractional
Brownian covariance},
arXiv:2602.19192 [math.PR], 2026.

\bibitem{Schoenberg}
I.~J.~Schoenberg,
\emph{Metric spaces and positive definite functions},
Trans.\ Amer.\ Math.\ Soc.\ \textbf{44} (1938), 522--536.

\bibitem{Bognar}
J.~Bogn\'ar,
\emph{Indefinite Inner Product Spaces},
Ergebnisse der Mathematik \textbf{78}, Springer, 1974.

\bibitem{BGL}
D.~Bakry, I.~Gentil, and M.~Ledoux,
\emph{Analysis and Geometry of Markov Diffusion Operators},
Grundlehren \textbf{348}, Springer, 2014.

\bibitem{BorceaBranden}
J.~Borcea and P.~Br\"and\'en,
\emph{P\'olya--Schur master theorems for circular domains and their boundaries},
Ann.\ of Math.\ (2) \textbf{170} (2009), 465--492.

\bibitem{RodgersTao}
B.~Rodgers and T.~Tao,
\emph{The de Bruijn--Newman constant is non-negative},
Forum Math.\ Pi \textbf{8} (2020), e6.

\bibitem{Hochberg}
K.~J.~Hochberg,
\emph{A signed measure on path space related to Wiener measure},
Ann.\ Probab.\ \textbf{6} (1978), 433--458.

\end{thebibliography}

\end{document}
